\begin{center}
{\large\bfseries A compact complex threefold fibred by tori over the projective line, and the six-sphere.}
\end{center}

\subsection*{Abstract}
We construct a compact connected complex manifold $X$ of dimension three together with a surjective holomorphic map $f\colon X \to \dbP^1$ whose fibres over the complement of three points $p_0, p_1, p_2$ are complex tori of dimension two, whose fibre over $p_0$ is a reduced irreducible rational surface with normal crossings (a del Pezzo surface of degree six with the opposite sides of its anticanonical hexagon identified), and whose fibres over $p_1$ and $p_2$ are multiple, of multiplicities $3$ and $4$, with bielliptic reductions $S_1, S_2$. The family over the thrice punctured sphere is given by explicit period functions on the upper half plane attached to the $(3,4,\infty)$ orbifold; the fibre over $p_0$ is filled in by a quotient of an infinite smooth toric threefold, and the fibres over $p_1, p_2$ by logarithmic transforms. We prove that $X$ is simply connected with the integral homology of $S^6$, hence diffeomorphic to $S^6$, and we compute $a(X) = 1$ with $f$ the algebraic reduction --- the N\'eron--Severi group of the very general fibre being infinite cyclic, generated by a class $\eta$ whose associated hermitian form $H_\eta$ has signature $(1,1)$ and which is therefore represented by no positive line bundle, so that the very general fibre has algebraic dimension $0$ --- together with $f_*\mathcal O_X = \mathcal O_B$, $R^1 f_* \mathcal O_X \cong \mathcal O_B \oplus \mathcal O_B(-1)$, $R^2 f_*\mathcal O_X \cong \mathcal O_B(-1)$, $h^{0,1} = 1$, $h^{0,2} = h^{0,3} = 0$, $h^{1,0} = h^{2,0} = 0$, $K_X \cong f^*\mathcal O_B(-1) \otimes \mathcal O_X(2S_2)$, $c_1 c_2 = 0$, $c_3 = 2$, $h^0(T_X) = 1$ and $\mathrm{Aut}^0(X) \cong \dbC^*$. The existence of such an $X$ is incompatible with \cite[Cor.\ 2.3]{CDP20}; the last section locates the point at which the two accounts diverge, by showing that $R^2 f_*(T_X \otimes L) \ne 0$ for every holomorphic line bundle $L$ on $X$, a consequence of the non-normality of the normal crossings fibre $W$.

\section{Introduction.}

\subsection{The result}

Throughout, $B := \dbP^1$ carries the affine coordinate $t$, and $p_1, p_2, p_0$ are the points $t=0$, $t=1$, $t=\infty$, with $t_c := 1/t$ the coordinate at $p_0$ and $B^\circ := B \setminus \Sigma$, $\Sigma := \{p_0,p_1,p_2\}$; further $m_1 := 3$, $m_2 := 4$. We write $\mathrm{dP}_6$ for the del Pezzo surface of degree $6$, the blow up of $\dbP^2$ in three non-collinear points, whose anticanonical cycle is a hexagon of six $(-1)$-curves. A bielliptic surface is a smooth compact complex surface $(E\times F)/G$ with $E,F$ elliptic curves and $G$ a finite group of translations of $E$ acting on $F$ with $F/G \cong \dbP^1$ \cite[V.5]{BHPV}, \cite{BdF}, \cite{Ser}. We write $\mathcal M(Y)$ for the field of meromorphic functions of a compact complex space $Y$, and $a(Y) = \mathrm{tr.deg}_{\dbC}\mathcal M(Y)$ for its algebraic dimension.

\begin{theorem}[Main Theorem]
There exist a compact connected complex manifold $X$ of dimension $3$ and a surjective holomorphic map $f\colon X \to B = \dbP^1$ with connected fibres, with the following properties.
\begin{enumerate}[label=(\arabic*)]
\item $f$ is a proper holomorphic submersion over $B^\circ$, and every fibre $f^{-1}(p)$ with $p \in B^\circ$ is a complex torus of dimension $2$.
\item The fibre $W := f^{-1}(p_0)$ is reduced, irreducible, and a normal crossings divisor in $X$: at every point of $W$ the function $t_c \circ f$ is, in suitable local coordinates, $z_1$, $z_1 z_2$ or $z_1 z_2 z_3$. The normalisation of $W$ is $\mathrm{dP}_6$, and $W$ is obtained from $\mathrm{dP}_6$ by identifying the three pairs of opposite sides of its hexagon of $(-1)$-curves. The double locus $D \subset W$ consists of three smooth rational curves, each passing through both triple points of $W$ and otherwise pairwise disjoint; and $e(W) = 2$.
\item The fibres over $p_1$ and $p_2$ are multiple: as divisors, $f^*(p_1) = 3S_1$ and $f^*(p_2) = 4S_2$, where $S_j = (f^{-1}(p_j))_{\mathrm{red}}$ is a smooth bielliptic surface, and the normal bundle $\mathcal O_X(S_j)|_{S_j}$ has order exactly $m_j$ in $\mathrm{Pic}(S_j)$.
\item $X$ is simply connected and $H_*(X;\dbZ) \cong H_*(S^6;\dbZ)$; in particular $b_2(X) = b_3(X) = 0$ and $e(X) = c_3(X) = 2 = e(W)$. Consequently $X$ is diffeomorphic to the $6$-sphere.
\item $a(X) = 1$, and $f$ is the algebraic reduction of $X$: $\mathcal M(X) = f^*\mathcal M(B) = f^*\dbC(t)$, that is, every meromorphic function on $X$ is the pull back under $f$ of a rational function on $B$; moreover the very general fibre of $f$ has algebraic dimension $0$.
\item $f_*\mathcal O_X = \mathcal O_B$, $R^1 f_*\mathcal O_X \cong \mathcal O_B \oplus \mathcal O_B(-1)$ and $R^2 f_*\mathcal O_X \cong \mathcal O_B(-1)$; moreover
\[
h^{0,1}(X) = 1, \qquad h^{0,2}(X) = h^{0,3}(X) = 0, \qquad h^{1,0}(X) = h^{2,0}(X) = h^{3,0}(X) = 0.
\]
In particular $\chi(\mathcal O_X) = 0$; $\mathrm{Pic}(X) = \mathrm{Pic}^0(X) \cong H^1(X,\mathcal O_X) \cong \dbC$, since $H^1(X;\dbZ) = H^2(X;\dbZ) = 0$; and the Fr\"olicher spectral sequence of $X$ does not degenerate at $E_1$, since $b_1(X) = 0 < h^{0,1}(X)$.
\item $K_X \cong f^*\mathcal O_B(-1) \otimes \mathcal O_X(2S_2)$; in particular $K_X$ is not a torsion element of $\mathrm{Pic}(X)$, because $\mathcal O_X(4S_2) = f^*\mathcal O_B(1)$ gives $K_X^{\otimes 4k} \cong f^*\mathcal O_B(-2k)$, whose space of sections is $H^0(B,\mathcal O_B(-2k)) = 0$ for $k \ge 1$ by $f_*\mathcal O_X = \mathcal O_B$, while $h^0(\mathcal O_X) = 1$ shows that the trivial bundle does have sections. The Chern classes satisfy $c_1(X) = c_2(X) = 0$ in $H^*(X;\dbZ)$, and the Chern numbers are $c_1 c_2 = 0$ and $c_3 = 2$, whence $\chi(X,T_X) = 1$ by Hirzebruch--Riemann--Roch.
\item $h^0(X,T_X) = 1$ and $\mathrm{Aut}^0(X) \cong \dbC^*$, acting along the fibres of $f$; the generating vector field is the vertical field $\xi$ determined by the monodromy invariant vector $\hat\delta$ of the period lattice, and the fixed locus of the action is a single double curve of $W$, a smooth rational curve (cf.\ \cite[Prop.\ 2.5]{HKP}).
\end{enumerate}
\end{theorem}

\begin{corollary}
The six-sphere $S^6$ carries an integrable complex structure, namely the one obtained by transporting the complex structure of $X$ along a diffeomorphism $X \to S^6$.
\end{corollary}

\begin{remark}
Since $b_2(X) = 0$, $X$ is not K\"ahler, is not Moishezon \cite[Lemma 1.4]{CDP98}, and does not belong to Fujiki's class $\mathcal C$. The invariants in items (6)--(8) agree with the known constraints on a compact complex threefold homeomorphic to $S^6$: $h^0(T_X) = 1 \le 2$ \cite[Cor.\ 1.3]{HKP}; $h^{0,1} = h^{0,2}+1$, $H^3(X,\mathcal O_X) = 0$ and $\mathrm{Pic}(X) = \mathrm{Pic}^0(X) \cong H^1(X,\mathcal O_X)$ \cite[Cor.\ 10.2]{HKP}; $K_X$ not torsion, as forced in \cite[\S10]{HKP}; and the fixed locus of the $\dbC^*$-action of item (8), computed in \S9, is of one of the two shapes permitted by \cite[Prop.\ 2.5]{HKP}.
\end{remark}

\begin{remark}[Hodge numbers not computed here]
All the Hodge numbers of $X$ are determined by item (6) except $h^{1,1}$ and $h^{1,2} = h^{2,1}$. Since $\sum_{p,q} (-1)^{p+q} h^{p,q} = e(X)$ on any compact complex manifold (the Euler characteristic is unchanged through the Fr\"olicher spectral sequence), equivalently by the Hirzebruch--Riemann--Roch identity $\sum_p (-1)^p \chi(\Omega^p_X) = c_3(X)$ (\S9.4), the values $\chi(\mathcal O_X) = 0$, $e(X) = 2$ and the vanishings of item (6) force $h^{1,1} = h^{2,2} = h^{1,2}+1$. Only $h^{1,2}$ is left undetermined; we expect $h^{1,2} = 1$ (whence $h^{1,1} = h^{2,2} = 2$). A computation would require the direct images $R^q f_*\Omega^1_X$ and $R^q f_*\Omega^2_X$, hence a description of the sheaves $\Omega^p_{X/B}$ along the three singular fibres, where they have torsion, together with an analysis of the corresponding Leray spectral sequences in the light of the non-degeneration recorded in item (6).

The construction is explicit. Over $B^\circ$ one has a family $J \to B^\circ$ of $2$-tori whose period matrix is written down in closed form in terms of three functions $\tau, \mu, \beta$ on the upper half plane $\hb_z$ uniformising the orbifold $(B;3,4,\infty)$: $\tau$ is the elliptic modular parameter, $j(\tau) = 1728\,t$, while $\mu$ and $\beta$ solve prescribed inhomogeneous transformation laws. The local monodromies at $p_1$ and $p_2$ have orders $3$ and $4$, and the local monodromy at $p_0$ is unipotent with square-zero logarithm. The three fibres over $\Sigma$ are filled in by three local models: at $p_0$ by a quotient of an infinite smooth toric threefold built on the $A_2$-triangulation of the plane, in the spirit of Mumford's construction of degenerating abelian varieties \cite{Mum72,AMRT,KKMS}; at $p_1$ and $p_2$ by logarithmic transforms in the sense of Kodaira \cite{Kod64}, with twist vectors $v_1 = \varepsilon$ and $v_2 = -\varepsilon'$, explicit vectors of the period lattice fixed by the respective local monodromies. The three twists are what makes $X$ simply connected: \S7 gives $\pi_1(X) \cong \dbZ/|p|\dbZ$ with $p = 12\ell_0 - 4\ell_1 - 3\ell_2$, where $(\ell_0,\ell_1,\ell_2)$ records the twist data, and the above choice yields $(\ell_0,\ell_1,\ell_2) = (0,1,-1)$, so $p = -1$.
\end{remark}

\subsection{Relation to \cite{CDP20}}

The Main Theorem is incompatible with a published result. Campana, Demailly and Peternell proved in \cite{CDP98} that a compact complex threefold $X$ with $b_2(X) = 0$ and $a(X) > 0$ has $c_3(X) = e(X) = 0$, and deduced that a compact complex threefold homeomorphic to $S^6$ has algebraic dimension $0$. That proof rests on \cite[Lemma 1.5]{CDP98}, which is incorrect as stated; the corrigendum \cite{CDP20} says so and re-proves the results in two cases. \cite[Thm.\ 2.1]{CDP20} treats the case in which there exists a non-constant meromorphic map $g\colon X \dashrightarrow \dbP^1$ which is not holomorphic (this case is written out in detail in \cite[Thm.\ 3.1, Cor.\ 3.2]{LRS}, following \cite{CDP98}), while \cite[Thm.\ 2.2]{CDP20} treats the case $H^1(X,\dbZ) = H^2(X,\dbZ) = 0$, $a(X) = 1$ with holomorphic algebraic reduction $f\colon X \to C$, concluding (a) $H^i(X,T_X\otimes M) = 0$ for $i \ne 1$ and generic $M \in \mathrm{Pic}^0(X)$, (b) $\chi(X,T_X\otimes M) \le 0$ and (c) $c_3(X) \le 0$; together these give \cite[Cor.\ 2.3]{CDP20}: if $X$ is homeomorphic to $S^6$ then $a(X) = 0$. The manifold of the Main Theorem has $H^1(X,\dbZ) = H^2(X,\dbZ) = 0$, $a(X) = 1$, holomorphic algebraic reduction $f\colon X \to \dbP^1$, and $c_3(X) = 2 > 0$, so that conclusion (c) fails; so does conclusion (b), since $c_1(X) = c_2(X) = 0$ and $\mathrm{Pic}(X) = \mathrm{Pic}^0(X)$ give, by Riemann--Roch, $\chi(X,T_X\otimes M) = \tfrac12 c_3(X) = 1 > 0$ for every $M \in \mathrm{Pic}(X)$. It therefore contradicts \cite[Thm.\ 2.2]{CDP20} and \cite[Cor.\ 2.3]{CDP20}, as well as \cite[Lemma 4.2]{CDP20} (the special case $g=0$, $b_1=0$ of \cite[Lemma 3.2]{CDP98}), \cite[Cor.\ 5.2]{CDP20} and \cite[Prop.\ 5.5]{CDP20} ($=$ \cite[Thm.\ 3.1, Rem.\ 3.8(1)]{CDP98}), whose conclusions fail for $X$ (their printed proofs presuppose trivial monodromy); this failure is repairable for the purpose of the proof of Cor.\ 2.3 at $X$ (\S10), unlike the reduction at the non-normal fibre, which is decisive. It does not contradict \cite[Thm.\ 2.1]{CDP20} or \cite{LRS}, whose hypothesis is the existence of some non-constant non-holomorphic meromorphic map $X \dashrightarrow \dbP^1$, and no such map exists on our $X$: by item (5) of the Main Theorem $\mathcal M(X) = f^*\dbC(t)$, so every non-constant meromorphic map $X \dashrightarrow \dbP^1$ is of the form $h \circ f$ with $h$ a non-constant rational function on $B$, hence is holomorphic.

Section 10 is devoted to this conflict and locates the step at which the argument of \cite{CDP20} and the example part company. The key point is the fibre $W = f^{-1}(p_0)$: it is a normal crossings divisor which is not normal, the six sides of its hexagon being glued to each other in three opposite pairs (item (2) of the Main Theorem). We prove (Theorem 10.5) that for every holomorphic line bundle $L$ on $X$ one has $R^2 f_*(T_X \otimes L) \ne 0$, the obstruction being supported at $p_0$ and arising from the section $df \otimes e$ of a sheaf supported on $W$ which exists precisely because $W$ is non-normal. The vanishing of this second direct image is the hypothesis under which \cite[Prop.\ 2.4]{CDP20} is invoked in the proof of \cite[Thm.\ 2.2]{CDP20}; for $X$ it fails, for every $L$, and so the proof of \cite[Thm.\ 2.2]{CDP20} does not cover $X$. Section 10 also shows, in the light of the monodromy group of $f$ computed in \S3, that the conclusions of \cite[Lemma 4.2]{CDP20}, \cite[Cor.\ 5.2]{CDP20} and \cite[Prop.\ 5.5]{CDP20} fail for $X$ as well. Section 10 discusses the two steps of the printed proofs that do not apply to $X$; the two statements cannot both be correct. For this reason the construction is given with every matrix, chart and identification explicit.

\subsection{Plan of the paper}

Section 2 fixes the linear algebra: the lattice $V \cong \dbZ^4$ with basis $(\gamma,u,w,\delta)$, the matrices $T_1, T_2$ of orders $3$ and $4$, the unipotent $T_0 = (T_1T_2)^{-1} = I + N$ with $N^2 = 0$, and the dual picture on $\Lambda = V^*$ with $A_j = A(T_j)$ and $M_0$. All facts used later are proved there: $\Lambda_{\mathrm{tor}} = \ker(M_0 - I) = \mathrm{im}(M_0-I)$, the unimodularity of $B_0\colon \bar\Lambda \to \Lambda_{\mathrm{tor}}$, the descriptions $\Lambda^{A_1} = \langle \varepsilon,\hat\delta\rangle$, $\Lambda^{A_2} = \langle \varepsilon',\hat\delta\rangle$, and the presentation $\Delta = \langle g_1,g_2 \mid g_1^3 = g_2^4 = 1\rangle \cong \dbZ/3 * \dbZ/4$ of the orbifold fundamental group of $(B;3,4,\infty)$, the $(3,4,\infty)$ triangle group, together with the representation $\rho_V\colon \Delta \to \mathrm{GL}(V)$, $\rho_V(g_j) = T_j$. The group $\Delta$ is not isomorphic to $\mathrm{PSL}_2(\dbZ) \cong \dbZ/2 * \dbZ/3$; there is a surjection $\Delta \to \mathrm{PSL}_2(\dbZ)$, and the modular parameter $\tau$ is equivariant along it.

Section 3 constructs the period family. Theorem 3.4 establishes existence and uniqueness, up to one additive constant $c_0$ which is then chosen, of $\tau,\mu,\beta$ on $\hb_z$ subject to the equivariance and boundedness conditions listed there, their behaviour at the two elliptic points and at the cusp, the equivariance $\Pi(gz) = R_g(z)\Pi(z)M_g$ of the period matrix $\Pi = [Z \mid I]$, and the nondegeneracy $\operatorname{Im}\tau \cdot D \ne 0$ which makes $L(z) = \Pi(z)\Lambda$ a lattice in $\dbC^2$ for every $z$.\footnote{The manifold $X$ depends on the choice of $c_0$; see Remark 6.4.} This yields the family $\mathcal T \to \hb_z$ of complex $2$-tori, the quotient $J \to B^\circ$, its marking by $\Lambda$, and its monodromy.

Section 4 constructs the filling at $p_0$. The fan $\mathfrak F$ of cones over the $A_2$-triangulation of $\dbR^2 \times \{1\}$ gives a smooth toric threefold $Y_{\mathfrak F}$, not of finite type, with a global holomorphic function $t_c$; the group $\Lambda \cong \dbZ^2$ acts by the maps $\Psi_{\bar\lambda} = (c_{\bar\lambda}(t_c),1)\cdot \Phi_{\bar\lambda}$. Theorem 4.5 shows that for small radius this action is free and properly discontinuous, that $N_0$ is a complex threefold with $f_0$ proper, that the fibres over $t_c \ne 0$ are the tori of $J$ under the tautological identification $E_0$, and that $W = f_0^{-1}(0)$ is as described in item (2). For the quotient to be a complex manifold proper over the disc, only $\det B_0 \ne 0$ and the holomorphy of $C(t_c)$ at $t_c = 0$ are used; it is the unimodularity $|\det B_0| = 1$ that makes the central fibre irreducible with normalisation a single $\mathrm{dP}_6$. No positivity or symmetry of $B_0$ enters. Toric generalities are taken from \cite{Ful93,Oda88,KKMS}.

Section 5 constructs the fillings at $p_1, p_2$ by logarithmic transforms. Theorem 5.4 shows that $\tilde g_j^{\log}$, which in flat coordinates reads $(z,x) \mapsto (g_j z, A_j x + v_j/m_j)$, generates a free $\dbZ/m_j$-action on $\mathcal T|_{\Delta_j}$ --- freeness holds because $3 \nmid \gamma(v_1)$ and $\gamma(v_2)$ is odd --- that the quotient $N_j$ is a complex threefold with $f_j^*(p_j) = m_j S_j$ and $S_j$ bielliptic, and that translation by the logarithmic section $\sigma_j$ conjugates $\tilde g_j^{\log}$ to $\tilde g_j^0$, identifying $N_j \setminus f_j^{-1}(p_j)$ with $J|_{D_j^*}$.

Section 6 glues $J, N_0, N_1, N_2$ and proves Theorem 6.2: $X$ is a compact connected complex threefold, $f$ is holomorphic and surjective, and $X$ is independent, up to biholomorphism over $B$, of the auxiliary choices (radii, discs, branches of logarithm).

Section 7 computes the topology. Using the $0$-section of $J$ to split $\pi_1(J|_{B^\circ}) \cong \Lambda \rtimes \pi_1(B^\circ)$, van Kampen gives $\pi_1(X) \cong \dbZ/|p|\dbZ$ with $p = 12\ell_0 - 4\ell_1 - 3\ell_2$, hence $\pi_1(X) = 1$ here (Theorem 7.17); Mayer--Vietoris, with the homology of $W$ and of the $S_j$ computed from the cell structures of the appendix, gives $H_*(X;\dbZ) \cong H_*(S^6;\dbZ)$ and $e(X) = 2$ (Theorem 7.22). The group $H_*(X;\dbZ)$ is in fact computed twice in that section: once along the Mayer--Vietoris route just described, which uses the retraction and collapse of \S7.2, and once along the Leray route of \S7.7 together with Appendix B, which, in combination with Theorem 7.17 and Corollary 4.8, makes no use of \S7.2; the recognition of $X$ as $S^6$ therefore does not depend on Proposition 7.2. For comparison, the threefold $X'$ obtained with $v_2 = +\varepsilon'$ has $\pi_1(X') \cong \dbZ/7\dbZ$.

Section 8 concludes that $X$ is diffeomorphic to $S^6$ (Theorem 8.1): $X$ is a closed smooth simply connected $6$-manifold with the integral homology of $S^6$, hence a homotopy $6$-sphere by Hurewicz and Whitehead \cite{Hat02}, hence homeomorphic to $S^6$ by \cite{Sma61}, and diffeomorphic to $S^6$ since $\Theta_6 = 0$ \cite{KM63}, with \cite{Mil65}.

Section 9 computes the invariants of Theorem 9.1. The algebraic dimension is determined by a direct computation on the fibres: from the explicit period matrix one finds that a class $E \in \Lambda^2 V$ lies in $\mathrm{NS}(F_z)$ if and only if an explicit linear form $P_E(z)$ in the five functions $1,\tau,\mu,\beta,6\mu^2-\tau\beta$ vanishes, and these five functions are linearly independent over $\dbC$; hence for $z$ outside a countable set one gets $\mathrm{NS}(F_z) = \dbZ\eta$ with $\eta = u\wedge w + 6\,\gamma\wedge\delta$, whose hermitian form $H_\eta$ has signature $(1,1)$, so that neither $\eta$ nor $-\eta$ is the class of a positive line bundle, while $\eta^2 = 12\cdot\mathrm{vol} \ne 0$ excludes $a(F_z) = 1$ as well; the very general fibre therefore has algebraic dimension $0$. Every meromorphic function on $X$ is therefore constant on the very general fibre, whence $\mathcal M(X) = f^*\mathcal M(B) = f^*\dbC(t)$ and $a(X) = 1$, with $f$ the algebraic reduction. Section 9 also gives a second proof, by exclusion: $a(X) \ge 1$ because $f^*$ embeds $\mathcal M(B)$ into $\mathcal M(X)$; $a(X) \ne 3$ because a Moishezon threefold has $b_2 > 0$ \cite[Lemma 1.4]{CDP98} while $b_2(X) = 0$ by \S7; and $a(X) \ne 2$, since otherwise $X$ carries a meromorphic map to a curve which is not holomorphic (as in the proof of \cite[Cor.\ 2.3]{CDP20}), whence $c_3(X) = 0$ by \cite[Thm.\ 2.1]{CDP20} (see also \cite[Thm.\ 3.1]{LRS}), contradicting $c_3(X) = e(X) = 2$. This route also yields the identification $\mathcal M(X) = f^*\mathcal M(B)$, as explained in Remark 9.9: once $a(X) = 1$, the field $\mathcal M(X)$ is algebraic over the subfield $f^*\mathcal M(B)$ of transcendence degree $1$, and $f^*\mathcal M(B)$ is algebraically closed in $\mathcal M(X)$ because $f$ has connected fibres \cite[\S3]{Ue75}; hence $\mathcal M(X) = f^*\mathcal M(B)$ and $f$ is the algebraic reduction. The section then computes the direct images $R^q f_*\mathcal O_X$ and the numbers $h^{0,q}$ of item (6), the vanishing of the holomorphic forms $h^{1,0} = h^{2,0} = h^{3,0} = 0$, the canonical bundle of item (7), the Chern numbers, and the vertical $\dbC^*$-action with $h^0(T_X) = 1$ and its fixed locus. Section 10 is the discussion of \S1.2; the appendix collects the matrices, the fan $\mathfrak F$, and the cell structures used in \S7.

\subsection{Context}

The question whether $S^6$ admits a complex structure goes back to Hopf \cite{Hopf48}; for its history and the state of the art we refer to the surveys \cite{AGr,Ang18}. Among the even-dimensional spheres only $S^2$ and $S^6$ carry almost complex structures; the structure on $S^6$ induced by octonionic multiplication is not integrable, its Nijenhuis tensor being nowhere zero, and the space of almost complex structures inducing a given orientation is connected, so that no homotopy-theoretic invariant can separate a hypothetical integrable structure from the octonionic one. The one general structural theorem in this direction is that of LeBrun \cite{LeB87}, anticipated by Blanchard: no complex structure on $S^6$ is orthogonal with respect to the round metric; the proof shows that such a structure would make $S^6$ K\"ahler, contradicting $b_2 = 0$, and the later refinements also use only $H^2(S^6) = 0$. The structure constructed here is not produced from a metric and carries no a priori relation to the round one, so \cite{LeB87} does not bear on it.

Compact complex threefolds with $b_2 = 0$ are classical: Calabi and Eckmann \cite{CE53} put complex structures on $S^3\times S^3$ and on products of odd-dimensional spheres; these, the Hopf threefolds $S^1\times S^5$, the quotients $\mathrm{SL}(2,\dbC)/\Gamma$ and the principal elliptic bundles of \cite[\S4]{LRS} furnish many examples, with $a = 0,1,2$ (see \cite[\S4]{CDP98}). What distinguishes $S^6$ among these is $e = 2 \ne 0$. The complex-analytic constraints on a hypothetical complex structure on $S^6$ have been studied in detail: Huckleberry, Kebekus and Peternell \cite{HKP} proved that such an $X$ is not almost homogeneous, that $h^0(T_X) \le 2$, $H^3(X,\mathcal O_X) = 0$, $h^{0,1} = h^{0,2}+1$ and $\mathrm{Pic}(X) = \mathrm{Pic}^0(X) \cong H^1(X,\mathcal O_X)$, and analysed the possible fixed-point sets of a $\dbC^*$-action; Lehn, Rollenske and Schinko \cite{LRS} obtained bounds on spaces of sections, and wrote out in detail the meromorphic-map case of \cite{CDP98,CDP20}, following \cite{CDP98}. Constructions of complex structures on $S^6$ have been claimed by Etesi \cite{Ete15} by entirely different means, from Samelson structures on $G_2$ and conjugate orbits; nothing here depends on them or bears on their validity.

The two mechanisms combined here are classical. Logarithmic transformations were introduced by Kodaira for elliptic surfaces \cite{Kod64} (see also \cite[V.13]{BHPV}), where they are the standard device for changing the fundamental group of an elliptic fibration without changing the base; the Dolgachev surface obtained from a rational elliptic surface by logarithmic transformations of coprime multiplicities is simply connected \cite{Dol81}, and the present construction is an analogue one dimension up, with $2$-tori in place of elliptic curves. Filling in a degenerating family of abelian varieties by a toroidal model over the nilpotent monodromy cone goes back to Mumford \cite{Mum72}, and in the arithmetic setting to \cite{AMRT}; the
\newpage
